% ch1_intro.tex — 第一章 绪论
\section{引言}

\subsection{研究背景}

在线评论是当前消费者行为研究的核心数据来源之一，但主流研究存在一个根本性的不对称：大量文献关注评论如何影响其他消费者的购买决策，却较少追问评论文本本身揭示了评论者怎样的心理结构，以及这些心理结构能否被系统提取和验证。这一缺口在方法论层面尤为突出——评论文本蕴含着丰富的信念、态度和意向信息，但现有研究缺乏将其转化为可量化心理变量的可靠工具。

Fishbein多属性态度理论\cite{fishbein1975}提供了一个成熟的理论框架：消费者对产品的总体态度（ATT）等于各属性信念强度（$b_i$）与信念评价（$e_i$）乘积的加总，即$\text{ATT} = \sum(b_i \times e_i)$。这一公式的核心价值在于将整体态度分解为可独立测量的信念成分，使研究者能够追踪态度形成的微观机制。然而，该理论在电商评论情境下的应用长期受限于测量手段：传统问卷无法大规模处理非结构化文本，词典方法难以准确捕捉否定结构和细粒度属性判断，监督学习方法则依赖大量人工标注数据\cite{grimmer2022}。

大语言模型（Large Language Models, LLMs）的出现为打破这一瓶颈提供了技术条件。LLM具备零样本推理能力，能够在无需专门标注数据的情况下，依据提示词中给定的理论框架执行信念提取和情感判断任务。这使得将Fishbein理论直接应用于大规模评论文本成为可能——将三个此前相互独立的要素连接起来：成熟的态度理论框架、海量的电商评论大数据、以及具备语义理解能力的LLM工具。然而，这一方法的测量效度——即LLM提取的多属性态度能否真正反映消费者的真实态度并预测其行为——仍是一个需要系统验证的实证问题。

\subsection{研究意义}

本研究的理论意义体现在方法论与理论应用两个层面。在方法论层面，本研究提出并系统验证了基于LLM的多属性态度提取框架，通过计算准确性验证、语义相似度检验和跨模型重测信度评估建立了一套可复用的质量控制标准，为"Text as Data"研究范式\cite{grimmer2022}在消费者心理变量提取中的应用提供了操作化依据。在理论层面，本研究将Fishbein多属性态度理论的测量逻辑从问卷情境迁移至非结构化评论文本，检验了多属性分解相对于单一评分的增量信息价值，并通过三数据源设计提供了一种在LLM文本分析研究中系统消除同源方差的方法论范式。

在实践层面，本研究为企业提供了低成本、可规模化的消费者态度监测工具。与传统问卷调查相比，LLM方法能够处理已积累的海量历史评论数据，在无需额外数据采集成本的前提下提取消费者对各产品属性的信念结构，支持数据驱动的产品改进和差异化营销策略。

\subsection{研究问题与研究目标}

\subsubsection{研究问题}

本研究聚焦于以下核心问题：

\begin{enumerate}
  \item LLM提取的多属性态度（ATT）能否有效预测消费者的客观评分行为（Rating）？若能，则证明基于Fishbein理论的LLM提取方法具有预测效度。
  \item 评论者经验（Reviewer Experience）是否调节ATT对Rating的预测强度？高经验评论者的多维评价框架是否会改变态度与评分之间的对应关系？
  \item 在控制Rating之后，ATT是否仍对行为意向（BI）具有显著增量预测力？若是，则证明多属性分解提供了单一评分所无法捕捉的信念结构信息。
\end{enumerate}

\subsubsection{研究目标}

基于上述研究问题，本研究旨在：

\begin{enumerate}
  \item 从在线评论中系统提取消费者信念（含情感评价$e_i$）和信念强度（$b_i$），计算多属性态度指数$\text{ATT} = \sum(b_i \times e_i)$；
  \item 通过四假设框架验证LLM提取方法的有效性：H1检验预测效度（$\text{ATT} \to \text{Rating}$），H2探讨调节效应，H3检验ATT对BI的预测力，H4检验增量效度（ATT在Rating之外对BI的额外预测力）；
  \item 探讨评论者经验对ATT--Rating关系的调节作用，检验消费者专业性对态度--行为一致性的影响；
  \item 建立LLM文本提取方法的可靠性验证标准，包括ATT计算准确性验证、语义相似度检验和重测信度评估。
\end{enumerate}

\subsection{研究内容}

本研究按以下结构展开，共分为十个章节：

\textbf{第一章\ 绪论}：阐述研究背景、研究意义、核心研究问题与目标，以及本文的整体研究内容框架。

\textbf{第二章\ 理论基础}：系统阐述本研究的核心理论框架，包括Fishbein多属性态度理论、计划行为理论、态度--行为一致性理论以及消费者专业性理论。本章为后续假设提出和结果解释提供理论支撑。

\textbf{第三章\ 文献综述}：围绕本研究的核心变量，分节梳理相关实证文献，涵盖：在线评论研究的两类视角、消费者信念与态度的实证研究、行为意向测量的实证进展、评论评分作为态度外显指标的文献依据、评论者经验与消费者专业性的实证研究，以及大语言模型在文本分析中的应用进展。

\textbf{第四章\ 研究假设}：基于第二章的理论基础和第三章的文献综述，提出本研究的四个核心假设，并对每个假设进行完整的逻辑推导：H1（多属性态度正向预测评分）、H2（评论者经验的调节效应，探索性）、H3（多属性态度正向预测行为意向）、H4（多属性态度提供超越评分的增量预测力）。

\textbf{第五章\ 研究设计}：说明研究方法选取理由（LLM方法对比传统方法）、研究设计框架、研究对象及代表性论证、数据来源与预处理、变量界定与操作化（含各变量的理论来源和测量方式）、LLM可靠性验证、数据分析方法、研究计划安排及当前进度。

\textbf{第六章\ 数据描述性统计}：报告样本基本特征、变量描述性统计、相关性分析等。

\textbf{第七章\ 研究结果}：报告LLM可靠性验证结果、H1/H2检验结果（有序Logistic回归及调节效应）和H3/H4检验结果（层次回归及增量效度）。

\textbf{第八章\ 讨论}：结合理论框架解释主要发现，讨论LLM方法的有效性、评论者经验的调节机制，以及各假设检验结果的理论意义。

\textbf{第九章\ 结论}：总结研究发现，阐述理论贡献与管理启示，讨论研究局限与未来展望。

\textbf{第十章\ 参考文献}：收录本研究引用的所有文献。

\textbf{第十一章\ 附录}：包含LLM提示词完整文本等补充材料。
